\documentclass[twocolumn, 10pt, a4paper]{article}

% Required Packages
\usepackage[utf8]{inputenc}
\usepackage{geometry}
\usepackage{graphicx}
\usepackage{amsmath, amssymb}
\usepackage{booktabs}
\usepackage{hyperref}
\usepackage{float}
\usepackage{caption}
\usepackage{subcaption}
\usepackage{enumitem}

% Margins for Journal Format
\geometry{
    a4paper,
    left=15mm,
    right=15mm,
    top=20mm,
    bottom=25mm,
}

% Title Data
\title{\textbf{Recommendation System: \\  Analysis on the Amazon Beauty Dataset}}
\author{Yunyu Guo, GuanMing Zhong \\ 
\small CS 6140: Machine Learning, Northeastern University}
\date{\today}

\begin{document}

\maketitle

\begin{abstract}
This project develops a recommendation system for the Amazon Beauty dataset, characterized by extreme sparsity and a long-tail distribution. We implement a ``Retrieve-then-Rank'' funnel architecture, utilizing a Two-Tower Deep Neural Network for candidate generation and comparing XGBoost against DeepFM for the ranking stage. Our evaluation reveals that while deep learning-based ranking (DeepFM) significantly outperforms gradient-boosted trees on implicit feedback data, the retrieval model struggles against simple collaborative filtering baselines (Item-KNN) due to insufficient interaction density. This report shows the data pipeline, feature engineering strategies based on metadata coverage, and an analysis of model performance under sparsity constraints.
\end{abstract}

\section{Introduction}
Modern recommendation systems must filter millions of items to find relevant products for users. The Amazon Beauty dataset presents a specific challenge: extreme data sparsity. With users averaging only 1.67 interactions, traditional collaborative filtering often fails due to the "Cold Start" problem.

To address this, we implemented a two-stage pipeline:
\begin{enumerate}
    \item \textbf{Retrieval:} Reducing the search space from 250k items to a candidate set of 200 using a dual-encoder (Two-Tower) model and Approximate Nearest Neighbors (ANN).
    \item \textbf{Ranking:} Re-ordering candidates using feature rich models (XGBoost and DeepFM) that leverage item metadata.
\end{enumerate}

\section{Data Analysis and Engineering}

\subsection{Dataset Statistics}
We utilize the \texttt{Amazon-Beauty} benchmark spanning 2010--2014.
\begin{itemize}[noitemsep]
    \item \textbf{Users / Items:} 1,210,271 / 249,274
    \item \textbf{Interactions:} 2,023,070
    \item \textbf{Sparsity:} 99.99\%
\end{itemize}
The interaction distribution follows a power law (Skewness: 28.6), meaning most users are ``inactive.''

\textbf{Data Preprocessing (5-Core Filter):} 
To mitigate this extreme sparsity, we applied a 5-core filter. This process iteratively removes users and items with fewer than 5 interactions until the dataset stabilizes. This ensures:
\begin{itemize}[noitemsep]
    \item \textbf{Learnability:} Every entity has sufficient interaction history for the model to detect patterns.
    \item \textbf{Valid Splitting:} Every user has enough data points to support a Train/Validation/Test split (preserving at least one item for testing).
\end{itemize}

\subsection{Data Cleaning and Feature Engineering}
To prepare the raw RecBole data for modeling, we implemented a cleaning pipeline using the Pandas library:

\begin{itemize}
    \item \textbf{Data Integrity:} We verified that there were zero duplicate user-item interaction pairs.
    \item \textbf{Missing Values:}
    \begin{itemize}
        \item \textit{Price (27\% missing):} We imputed missing values using the median price of the training set to be robust against outliers.
        \item \textit{Brand (51\% missing):} Due to low coverage, we treated missing brands as a distinct category, imputing them with an ``Unknown'' token.
    \end{itemize}
    \item \textbf{Outlier Handling \& Normalization:} Price and Sales Rank distributions were highly right-skewed. Instead of clipping outliers (which represent valid high-value or popular items), we applied a $\log(1+x)$ transformation to normalize these features for the XGBoost and DeepFM models.
    \item \textbf{Feature Engineering:} We parsed the raw JSON-like \texttt{categories} string to extract the \textit{Primary Category} (root node), simplifying the hierarchical tree into a usable categorical feature. Timestamps were normalized to [0,1] based on the training interval.
\end{itemize}

\section{Methodology}

\subsection{Retrieval Stage: Two-Tower}
We implemented a Dual Encoder architecture \cite{sampling_bias} where users and items are mapped to a shared 128-dimensional embedding space.
\begin{itemize}[noitemsep]
    \item \textbf{Objective:} Binary Cross Entropy (BCE) with a 1:4 positive-to-negative sampling ratio.
    \item \textbf{Indexing:} We utilized FAISS (IndexFlatIP) to perform efficient similarity search over the 250k item vectors.
    \item \textbf{Baselines:} We compared this against \textit{Most Popular} and \textit{Item-KNN} (Sparse Matrix Cosine Similarity).
\end{itemize}

\subsection{Ranking Stage: XGBoost vs DeepFM}
The ranking models score the top-N candidates. To combat the high dimensionality and sparsity of the input space, we applied distinct regularization strategies:

\begin{itemize}
    \item \textbf{XGBoost (Tree-based) \cite{xgboost}:} We utilized Early Stopping (patience=30 rounds) to halt training when validation loss plateaued, preventing the trees from modeling noise. We also applied L2 regularization ($\lambda=1.0$) to leaf weights to penalize overly complex trees.
    \item \textbf{DeepFM (Neural) \cite{deepfm}:} Combines Factorization Machines \cite{rendle_fm} with a Deep Neural Network. We applied a Dropout Rate of 0.2 to the deep layers to prevent neuron co-adaptation, ensuring the model learned robust distributed representations rather than relying on specific weights.
\end{itemize}


\section{Experiments and Results}

\subsection{Evaluation Protocol}
We utilized RecBole's time ordered (TO) split strategy with a ratio of 80/10/10 for train/validation/test.

\textbf{Note on Cross-Validation:} While K-Fold Cross-Validation is standard for tabular data, we prioritized a \textbf{Time-Ordered Holdout Strategy} over K-Fold CV. Standard random cross-validation in recommendation systems introduces look ahead bias (temporal leakage) by allowing the model to train on future user actions to predict past ones. The validation set was strictly used for hyperparameter tuning and early stopping.

For consistent comparison across models, we implemented a standardized random-negative evaluation:
\begin{itemize}
    \item Sample 10,000 users from the test set
    \item For each user: 1 positive item + 5 random negatives 
    \item Compute ranking metrics on this 6-item list per user
\end{itemize}

\subsection{Evaluation Metrics}
We selected metrics aligned with the specific objectives of each stage:
\begin{itemize}
    \item \textbf{Retrieval (Recall@K, Hit Rate):} Since the retrieval stage aims to capture relevant items regardless of their specific order within the candidate set, we utilize unweighted metrics like Recall@K to measure the proportion of relevant items successfully recalled from the full catalog.
    \item \textbf{Ranking (NDCG@K, AUC):} The ranking stage must prioritize relevant items at the top of the list. We use \textbf{NDCG@K} (Normalized Discounted Cumulative Gain) as our primary metric because it applies a position-based discount, heavily penalizing relevant items appearing lower in the recommendation list. \textbf{AUC} (Area Under the Curve) serves as a diagnostic metric for measuring the model's ability to discriminate between positive and negative samples in a pairwise manner.
\end{itemize}

\subsection{Retrieval Performance}
We evaluated retrieval using Hit Rate (HR) on the full catalog.

\begin{table}[h]
\centering
\begin{tabular}{lccc}
\toprule
\textbf{Model} & \textbf{HR@10} & \textbf{HR@20} & \textbf{HR@50} \\
\midrule
Most Popular & 0.77\% & 1.32\% & 2.41\% \\
\textbf{Item-KNN} & \textbf{1.23\%} & \textbf{1.80\%} & \textbf{2.94\%} \\
Two-Tower (Neural) & 0.72\% & 1.11\% & 2.01\% \\
\bottomrule
\end{tabular}
\caption{Retrieval performance on full catalog search.}
\label{tab:retrieval}
\end{table}

Item-KNN outperformed the Neural Two-Tower model by 41\% at HR@10. In extremely sparse datasets (1.7 clicks/user), simple co-occurrence statistics often provide a stronger signal than learned embeddings, which struggle to converge without sufficient data density.

\subsection{Ranking Performance}
Ranking was evaluated using a "1 Positive + 5 Negatives" protocol per user ($N=10,000$ test users). Note: Recall@10 is 1.0 by definition in this protocol (6 items total); NDCG is the discriminative metric.

\begin{table}[h]
\centering
\begin{tabular}{lcc}
\toprule
\textbf{Model} & \textbf{NDCG@10} & \textbf{NDCG@20} \\
\midrule
XGBoost (Content Only) & 0.6120 & 0.6120 \\
XGBoost (With IDs) & 0.6436 & 0.6436 \\
\textbf{DeepFM} & \textbf{0.9706} & \textbf{0.9706} \\
\bottomrule
\end{tabular}
\caption{Ranking performance (Random Negative Protocol).}
\label{tab:ranking}
\end{table}

\section{Discussion}

\subsection{The Dominance of Embeddings in Ranking}
DeepFM achieved near perfect ranking (NDCG $\approx$ 0.97), placing the true positive item at Rank 1 for 97\% of users. In contrast, XGBoost (NDCG $\approx$ 0.64) placed the positive item at Rank 1 for only 32.2\% of users.

This disparity highlights that \textbf{Collaborative Filtering signals (IDs)} are far more predictive than \textbf{Content signals (Metadata)} for this dataset. DeepFM learns an interaction function between UserID and ItemID. XGBoost, while powerful on tabular data, struggles to learn these high-cardinality interactions without explicit embedding layers.

\subsection{Overfitting Analysis: Content vs. IDs}
Our XGBoost experiments demonstrated the Bias-Variance tradeoff. The model using User/Item IDs exhibited high variance (overfitting), with a significant generalization gap (Train AUC 0.92 vs. Test 0.70). The high cardinality of IDs allows the trees to ``memorize'' specific user history.

Conversely, the ``Content-Only'' ablation (removing IDs) reduced this gap to just 0.05 (Train 0.71 vs. Test 0.66), showing better stability but higher bias (underfitting), as it lacked the capacity to capture personalized user preferences.


\section{Conclusion}
We implemented a full recommendation pipeline. Our analysis suggests that for highly sparse datasets like Amazon Beauty:
1.  Retrieval: Simple memory based methods (KNN) outperform complex neural encoders (Two-Tower) due to data scarcity.
2.  Ranking: Deep learning models (DeepFM) vastly outperform gradient boosting (XGBoost) by effectively learning latent representations for high-cardinality ID features.

Future work should focus on enriching the Two-Tower model with the metadata features identified in our EDA to help bridge the sparsity gap in the retrieval stage.

% References
\begin{thebibliography}{9}

\bibitem{recbole}
Zhao, W. X., et al. (2021). RecBole: Towards a Unified, Comprehensive and Efficient Framework for Recommendation Algorithms. \textit{CIKM 2021}.

\bibitem{deepfm}
Guo, H., Tang, R., Ye, Y., Li, Z., \& He, X. (2017). DeepFM: A Factorization-Machine based Neural Network for CTR Prediction. \textit{IJCAI 2017}.

\bibitem{sampling_bias}
Yi, X., Yang, J., Hong, L., et al. (2019). Sampling-Bias-Corrected Neural Modeling for Large Corpus Item Recommendations. \textit{RecSys 2019}.

\bibitem{xgboost}
Chen, T., \& Guestrin, C. (2016). XGBoost: A Scalable Tree Boosting System. \textit{KDD 2016}.

\bibitem{rendle_fm}
Rendle, S. (2010). Factorization Machines. \textit{ICDM 2010}.

\bibitem{youtube_dnn}
Covington, P., Adams, J., \& Sargin, E. (2016). Deep Neural Networks for YouTube Recommendations. \textit{RecSys 2016}.

\end{thebibliography}

\end{document}